\begin{table}[htbp]
\centering
\caption{Standardized Effect Sizes}
\label{tab:sde}
\begin{threeparttable}
\begin{tabular}{lcccccc}
\toprule
Outcome & $\hat{\beta}$ & SE & SD($Y$) & SDE & SE(SDE) & Classification \\
\midrule
Separation Rate & -0.0086 & 0.0048 & 0.1175 & -0.073 & 0.041 & Moderate negative \\
New Hire Rate & -0.0109 & 0.0055 & 0.1430 & -0.077 & 0.039 & Moderate negative \\
Log Earnings & 0.0074 & 0.0094 & 0.3647 & 0.020 & 0.026 & Small positive \\
Log Employment & 0.0307 & 0.0142 & 1.8677 & 0.016 & 0.008 & Small positive \\
\bottomrule
\end{tabular}
\begin{tablenotes}[flushleft]\small
\item \textit{Notes:} \textbf{Country:} United States. \textbf{Research question:} Do predictive scheduling (``fair workweek'') laws, which mandate advance notice of work schedules in food services and retail, reduce worker separations and new hiring in covered industries? \textbf{Policy mechanism:} Laws require employers with 500+ employees in retail and food services to provide 14-day advance schedule notice, pay ``predictability pay'' for last-minute changes, and offer additional hours to existing workers before hiring new ones; this increases the cost of flexible scheduling and creates incentives to stabilize the incumbent workforce. \textbf{Outcome definition:} Separation rate (quarterly separations divided by beginning-of-quarter employment), new hire rate (new hires divided by employment), log average monthly earnings for stable jobs, and log beginning-of-quarter employment; all from Census QWI. \textbf{Treatment:} Binary; county-level adoption of a predictive scheduling ordinance (6 cities and 1 state, staggered 2015--2020). \textbf{Data:} Census Quarterly Workforce Indicators (QWI/LEHD), county $\times$ NAICS sector $\times$ quarter, 2013--2019, all U.S. counties with non-missing data ($N = 326,358$). \textbf{Method:} Triple-difference (treated county $\times$ covered industry $\times$ post) with county$\times$industry, industry$\times$quarter, and county$\times$quarter fixed effects; standard errors clustered at the county level. \textbf{Sample:} All U.S. counties with positive employment in food services (NAICS 72), retail (NAICS 44-45), manufacturing (NAICS 31-33), and professional services (NAICS 54); restricted to 2013Q1--2019Q4 to avoid COVID contamination; 45 treated counties across 5 adoption cohorts. SDE $= \hat{\beta} / \text{SD}(Y)$ where SD($Y$) is the pre-treatment standard deviation of the outcome among treated-industry observations. Classification refers to magnitude, not statistical significance: Large ($|$SDE$| > 0.15$), Moderate ($0.05$--$0.15$), Small ($0.005$--$0.05$), Null ($< 0.005$).
\end{tablenotes}
\end{threeparttable}
\end{table}
